% =============================================================================
%  BibHallu-Bench -- BGK601 Final Raporu
%  Yazar: Novruz Amirov (707241004) | ITU Siber Guvenlik Muh. ve Kriptografi
%  Derleme: latexmk -pdf main      (veya: pdflatex; biber; pdflatex; pdflatex)
% =============================================================================
\documentclass[11pt,a4paper]{article}

% --- Dil ve karakter kodlamasi ---
\usepackage[utf8]{inputenc}
\usepackage[T1]{fontenc}
\usepackage[turkish]{babel}
\usepackage{amsmath}
\usepackage{amssymb}

% --- Sayfa duzeni ve tipografi ---
\usepackage[a4paper,margin=2.5cm]{geometry}
\usepackage{setspace}
\setstretch{1.25}                    % proje doc: 1.15-1.5 araliginda
\usepackage{microtype}
\usepackage{parskip}
\usepackage{titlesec}

% --- Tablo, grafik, link ---
\usepackage{graphicx}
\usepackage{booktabs}
\usepackage{multirow}
\usepackage{tabularx}
\usepackage{array}
\usepackage{caption}
\usepackage{subcaption}
\usepackage{xcolor}
\usepackage{enumitem}
\usepackage{listings}
\usepackage[hidelinks]{hyperref}
\usepackage{cleveref}

% --- Kaynakca (IEEE stili) ---
\usepackage[backend=biber,style=ieee,sorting=none]{biblatex}
\addbibresource{refs.bib}

% --- Placeholder makrolari (Stage 3 sonra doldurulacak yerler) ---
\newcommand{\TBD}{\textcolor{red}{\textbf{[?]}}}
\newcommand{\todoresult}[1]{%
  \par\medskip\noindent\fcolorbox{red!50}{red!8}{%
    \begin{minipage}{0.96\linewidth}\small
      \textcolor{red}{\textbf{[STAGE-3 SONUCU BEKLENIYOR]}} #1
    \end{minipage}}\par\medskip}
\newcommand{\figplaceholder}[2]{%
  \fbox{\begin{minipage}{0.9\linewidth}\centering\vspace{1cm}%
    \textcolor{red}{\textbf{[Sekil pending: #1]}}\\[0.2em]%
    \texttt{#2}\vspace{1cm}\end{minipage}}}

\lstset{basicstyle=\ttfamily\footnotesize, breaklines=true,
        frame=single, framesep=4pt, columns=fullflexible}

% --- Baslik bilgileri ---
\title{\textbf{Arastirma Makalelerinde Bibliyografya
       Halusinasyonlarinin LLM Tabanli Tespiti:\\
       Cok Modelli Karsilastirmali Degerlendirme ve RAG Ablasyonu}\\[0.6em]
       \large BGK601 -- Final Raporu}
\author{Novruz Amirov (707241004)\\
        Istanbul Teknik Universitesi\\
        Siber Guvenlik Muhendisligi ve Kriptografi Ana Bilim Dali\\
        Ders Koordinatoru: Prof. Dr. Kemal Bicakci}
\date{31 Mayıs 2026 \quad ---\quad Bahar 2025--2026}

% =============================================================================
\begin{document}
\maketitle
\thispagestyle{empty}

\begin{center}\small
\fbox{\begin{minipage}{0.85\linewidth}\small\textbf{Yapay Zeka Araçları Kullanım Notu.}
Bu raporun ve sunum slaytlarının hazırlanmasında Claude (Anthropic) ve
ChatGPT (OpenAI) araçlarından kod iskeleti, metin düzenleme ve kaynak
tarama desteğinde yararlanılmıştır. Tüm deneysel sonuçlar, benchmark
anotasyonu ve nihai metin öğrenci tarafından denetlenmiştir; YZ çıktıları
ground-truth olarak kullanılmamıştır (bkz.~Ek~\ref{appx:ai}).\end{minipage}}
\end{center}

% --- Ozet (1 sayfa) -----------------------------------------------------------
\section*{Ozet (Executive Summary)}
Buyuk dil modellerinin (LLM) akademik metin uretiminde yayginlasmasiyla
birlikte, bu modellerin urettigi bibliyografya bolumlerinde var olmayan veya
hatali atiflanmis kaynaklarin (bibliyografya halusinasyonlari) gorulme sikligi
ciddi bir akademik butunluk sorunu olusturmustur. Bu calismada, akademik
referanslarda bes farkli halusinasyon turunu (H1: tamamen uydurma referans,
H2: metadata halusinasyonu, H3: semantik uyumsuzluk, H4: kronolojik
halusinasyon, H5: dogru referans/negatif) ayirt eden 300 ornekli ozgun bir
benchmark veri kumesi tasarlanmis ve dort LLM uzerinde sistematik bir
karsilastirma yurutulmustur: iki kapali kaynak API modeli (GPT-4o,
GPT-4o-mini) ile iki acik kaynak yerel model (Llama~3.1-8B, Qwen2.5-7B,
Ollama uzerinde Apple~M2~Pro/16~GB birlesik bellek). Acik kaynak modellere
ayrica CrossRef tabanli bir RAG (retrieval-augmented generation) katmani
uygulanmis; istatistiksel anlamlilik bootstrap guven araliklari ve McNemar
testleriyle raporlanmistir.

Ana bulgular: ikili halusinasyon tespiti F1 skoru GPT-4o icin
\textbf{0.911} (en yuksek), Qwen2.5-7B icin \textbf{0.071} (en dusuk) olarak
olculmustur. RAG entegrasyonu, acik kaynak modellerde ortalama
$\Delta\text{F1}=+0.320$ kazanim saglamis (Llama: $+0.231$, Qwen: $+0.409$);
RAG'li Llama'nin F1'i \textbf{0.792}'ye yukselerek GPT-4o-mini'ye (0.862)
\emph{0.07 puanlik} bir farka yaklasmistir. Tum modellerde JSON ayikla-
bilirligi \%100, calistirmalar arasi tutarlilik (Fleiss~$\kappa$) 0.90 ile
1.00 araligindadir. Maliyet-basarim ekseninde GPT-4o-mini
(\$0.0003/dogru karar) en verimli kapali kaynak; acik kaynak modeller sifir
marjinal maliyetle yerel olarak calistirilabilmektedir. Tum pairwise farklar
McNemar testiyle istatistiksel olarak anlamlidir ($p < 0.01$).

\textbf{Anahtar Kelimeler:} LLM degerlendirmesi, halusinasyon tespiti,
bibliyografya, RAG, akademik butunluk, CrossRef, ChromaDB.

\newpage

% --- 1. Giris ---------------------------------------------------------------
\section{Giris ve Motivasyon}
\label{sec:giris}

Buyuk dil modelleri (Large Language Models, LLM), arastirma destek
sureclerinin --- yazim, ozetleme, kaynak tarama --- ayrilmaz bir bileseni hâline
gelmistir. Vaswani ve ark.\ tarafindan onerilen Transformer mimarisi
\cite{vaswani2017attention} ve onu izleyen olcekli on-egitim paradigmasi
\cite{devlin2019bert,brown2020gpt3}, bu modellerin dogal dil uretimindeki basarisini
beklentilerin uzerine cikarmis; ancak bu basarinin yaninda \emph{halusinasyon}
adi verilen sistematik bir hata sinifi da gun yuzune cikmistir. Ji ve ark.\
halusinasyonu, ``modelin kendi parametrik bilgisinden veya verili kaynaktan
sapan, dogrulanabilir bicimde yanlis icerik uretmesi'' olarak tanimlar
\cite{ji2023survey}.

\paragraph{Bibliyografyaya ozgu sorun.}
Halusinasyon, ozellikle akademik atiflarda kritik bir bicim alir. Walters ve
Wilder, ChatGPT'nin urettigi bibliyografyalarda yer alan kaynaklarin yarisindan
fazlasinin ya hic var olmadigini ya da en az bir alaninda hatali oldugunu
raporlamistir \cite{walters2023fabrication}. Saglik bilimleri alaninda yapilan
benzer calismalar da, tibbi metin uretiminde uretilen referanslarin %47-%73
oraninda fabrikasyon icerdigini gostermistir
\cite{bhattacharyya2023high,athaluri2023exploring}. Bu bulgular yalnizca
akademik usulsuzlugu degil, ayni zamanda yapilan iddialarin
\emph{izlenebilirligini} de tehdit eder.

\paragraph{Mevcut yaklasimlarin yetersizligi.}
Bibliyografya dogrulamasinda kullanilan klasik yontemler iki uca yerlesir:
(i) tamamen elle dogrulama, ki olceklenmesi mumkun degildir; (ii) basit
DOI/URL kontrolu, ki yalnizca \emph{yokluk}u tespit edebilir, anlam diizeyindeki
celiskileri (yanlis dergi, baglama uymayan icerik, imkansiz yil) yakalayamaz.
GPTZero'nun Hallucination Detector gibi yeni nesil ticari araclar iddia
duzeyinde kaynak onerisi sunmakta; ancak \emph{var olan} bir referansin icerik
ve baglam uyumunu analiz edememektedir.

\paragraph{Bu calismanin katkilari.}
Bu rapor, LLM'leri hem \emph{degerlendirme aracilari} hem de
\emph{degerlendirilen ozne} olarak konumlandirir ve su katkilari sunar:

\begin{enumerate}[nosep]
\item Bes halusinasyon turunu (H1--H5) ayirt eden, programatik insa + insan
dogrulamasiyla olusturulmus, 300 ornekli ozgun bir \emph{benchmark veri
kumesi}.
\item Yapilanmis JSON ciktisi ve cogunluk karariyla calisan, $\geq$3 bagimsiz
calistirmali, model-agnostik bir \emph{degerlendirme cercevesi};
bootstrap guven araliklari ve McNemar testleriyle istatistiksel anlamliligi
ile birlikte.
\item Iki kapali kaynak (GPT-4o, GPT-4o-mini) ve iki acik kaynak (Llama~3.1-8B,
Qwen2.5-7B) modelin esit kosullarda \emph{karsilastirmali analizi}.
\item CrossRef tabanli bir RAG katmaninin acik kaynak modellere getirdigi
kazanimin \emph{ablasyon olcumu}: parametric memory--only --- {} RAG farki
$\Delta$F1, $\Delta$Precision, $\Delta$Recall biciminde.
\item Tum kod, prompt sablonu ve veri kartinin acik kaynak olarak
yayinlanmasi (MIT lisansli GitHub deposu).
\end{enumerate}

\paragraph{Rapor yapisi.}
Bolum~\ref{sec:related} ilgili calismalari ozetler; Bolum~\ref{sec:yontem}
yontem-bilimi sunar; Bolum~\ref{sec:sonuc} deneysel bulgulari verir;
Bolum~\ref{sec:tartisma} tartisma ve sinirlamalari; Bolum~\ref{sec:gelecek}
sonuc ve gelecek calismalari kapsar.

% --- 2. Ilgili Calismalar ---------------------------------------------------
\section{Ilgili Calismalar}
\label{sec:related}

\subsection{LLM halusinasyonu: taksonomi ve olcum}
Ji ve ark.\ \cite{ji2023survey} halusinasyonu \emph{intrinsic} (verili
girdiden sapma) ve \emph{extrinsic} (verili girdide bulunmayan ama dogrulanmayan
bilgi) olmak uzere ikiye ayirir. Huang ve ark.\ \cite{huang2023hallucination}
bunu daha ince ayirimlara goturur: \emph{faktualite halusinasyonu},
\emph{tutarlilik halusinasyonu} ve \emph{ihtimal halusinasyonu}. Bibliyografya
halusinasyonu, tipik olarak faktualite halusinasyonunun bir alt sinifidir.
Olcum tarafinda, Maynez ve ark.\ \cite{maynez2020faithfulness} insan
degerlendirici tabanli sadakat olculerini, Manakul ve ark.\
\cite{manakul2023selfcheckgpt} kara kutu modelleri icin \emph{ozkontrol}
yaklasimini onermistir. Min ve ark.\ FActScore'da \cite{min2023factscore} uzun
ciktilarin atomik gercek ifadelere ayrilarak skorlanmasini onerir; bu fikir,
bizim referans-basina-karar yapimizla yapisal olarak benzerlik gosterir.

\subsection{Bibliyografya/atif halusinasyonu}
Walters ve Wilder \cite{walters2023fabrication} ChatGPT'nin urettigi
84 bibliyografya icin elle dogrulama yapmis; referanslarin %55'inde en az bir
alan hatasi (yanlis yazar, yil, sayfa numarasi vb.), %18'inde ise referansin
tamamen uydurma oldugunu raporlamistir. Tibbi metinlerde
Bhattacharyya ve ark.\ \cite{bhattacharyya2023high} %47, Athaluri ve ark.\
\cite{athaluri2023exploring} %69 fabrikasyon orani bildirir. Bu sayisal
farkliliklar, hem alan baginliligini hem de \emph{tek tip benchmark
eksikligini} gosterir; bizim calismamiz bu eksikligi bes kategorili bir
taksonomiyle gidermeye yoneliktir.

\subsection{Retrieval-Augmented Generation (RAG)}
Lewis ve ark.\ \cite{lewis2020rag} dis bilgi tabaninin parametric memory'nin
sinirlarini astigi senaryolari sistemlestirmis; Borgeaud ve ark.\ RETRO
mimarisiyle \cite{borgeaud2022retro} milyarlarca token icinden retrieval ile
modelin daha kucuk parametrik agirliklarla daha iyi performans gosterebilecegini
gostermistir. Izacard ve Grave \cite{izacard2021leveraging} retrieval'in
acik alan QA'da getirdigi kazanimi belgeler. Yakin zamandaki kapsamli RAG
literatur taramasi Gao ve ark.\ \cite{gao2024ragsurvey} tarafindan
yapilmistir. Bibliyografya dogrulamasi, ozellikle ``referansin gercekligi'' icin
RAG'in dogal bir uygulama alanidir: bilgi tabanindaki en yakin kayitla yapilan
karsilastirma, modelin uydurma sezgisini disardan deneyimle besler.

\subsection{Embedding ve retrieval altyapilari}
Reimers ve Gurevych'in \cite{reimers2019sbert} Sentence-BERT'i, semantik
benzerlik gorevlerinde cumle/paragraf duzeyinde gomulmeleri sistemlestirmistir;
bu calisma onlarin \texttt{all-MiniLM-L6-v2} modelini retriever embedder olarak
kullanir.

\subsection{Prompt muhendisligi}
Chain-of-thought \cite{wei2022cot} ve self-consistency \cite{wang2023selfconsistency}
teknikleri akil yurutme gerektiren gorevlerde tutarli kazanim saglar; H3
(semantik uyumsuzluk) gibi kategorilerimizde bu teknikleri ablasyon olarak
denedik (bkz.~Bolum~\ref{sec:tartisma}).

\subsection{Siber guvenlik baglaminda LLM}
Bhatt ve ark.\ CyberSecEval \cite{bhatt2023cyberseceval} LLM'lerin guvenlik
tartismalarinda ne kadar guvenli ve dogru oldugunu sorgular; Yao ve ark.\
\cite{yao2024survey} ise LLM'lerin guvenlik etkilerini iki yonlu ele alir.
Akademik butunlukteki halusinasyon, dolayli olarak dezenformasyon ve sosyal
muhendislik vektorlerine baglandigi icin siber guvenlik calismalari kapsaminda
mesrudur.

\subsection{LLM degerlendirme metodolojisi}
Liang ve ark.\ HELM cerceveiyle \cite{liang2023helm} cok boyutlu LLM
degerlendirmesinin (dogruluk + sertlik + maliyet) gerekliligini ortaya koymus;
Zheng ve ark.\ ``LLM-as-a-Judge'' yaklasiminin sinirlarini gosterir
\cite{zheng2023judging}. Bu calisma, otomatik metriklerle (F1, Fleiss~$\kappa$,
bootstrap guven araliklari) insan dogrulamasini birlikte kullanir.

% --- 3. Yontem --------------------------------------------------------------
\section{Yontem}
\label{sec:yontem}

\subsection{Uygulama tanimi ve kapsam}
Calismanin uygulama alani, akademik makalelerin \emph{bibliyografya bolumlerinde}
yer alan referanslarin LLM tabanli analizidir. Sistem, IEEE/ACM/APA biciminde
yazilmis Ingilizce bir referansi (gerekirse atif baglamiyla birlikte) girdi
olarak alir; cikti olarak \texttt{\{hallucination, type, confidence, reason\}}
alanlari iceren bir JSON karari uretir. Hedef kullanicilar arastirmacilar,
editorler ve akademik etik birimlerdir. Kapsam disinda kalan unsurlar: tam
metin dogrulamasi, in-text citation analizi, Ingilizce disindaki dillerde
atiflar ve patent/teknik rapor referanslari.

\subsection{Halusinasyon turleri ve kategorizasyon}
\label{sub:cat}
Calismadaki bes halusinasyon turunu Tablo~\ref{tab:cat} ozetler. Bu sema, hem
mevcut literaturdeki sezgisel siniflandirmalari
\cite{walters2023fabrication,huang2023hallucination} hem de bibliyografyalara
ozgu hata oruntulerini hesaba katar.

\begin{table}[t]
\caption{Halusinasyon kategorileri (H1--H5) ve insa mantigi.}
\label{tab:cat}\small
\begin{tabularx}{\columnwidth}{l l X}
\toprule
\textbf{Kod} & \textbf{Kategori} & \textbf{Insa mantigi (deterministik)} \\
\midrule
H1 & Var olmayan referans & Olusturma kurallariyla uretilen, hicbir veritabaninda yer almayan tamamen sentetik atif. \\
H2 & Metadata halusinasyonu & Gercek bir kaydin yazar/yil/dergi alaninin biri rastgele bozulur. \\
H3 & Semantik uyumsuzluk & Gercek bir kayit, konusuyla ilgisiz bir atif baglamina yerlestirilir. \\
H4 & Kronolojik halusinasyon & Yil, gelecek tarihe veya teknolojinin var olusundan oncesine cekilir. \\
H5 & Dogru referans (negatif) & Tum alanlari dogrulanmis gercek kayit. \\
\bottomrule
\end{tabularx}
\end{table}

\subsection{Gereksinim analizi ve kriter cercevesi}
\label{sub:kriter}
LLM'nin karsilamasi gereken fonksiyonel gereksinimler: (i)~yapilanmis JSON
ciktisi, (ii)~$\geq$8K tokenli baglam penceresi (20--60 referansli
bibliyografyalar icin), (iii)~CrossRef/OpenAlex tarzi metadata semasini
yorumlayabilme, (iv)~cok adimli akil yurutme (retrieval $\to$ verification
$\to$ scoring). Fonksiyonel olmayan gereksinimler: gizlilik (kurumsal dagitim
icin on-premise tercih), gecikme (batch toleransli, referans basina $\leq$30~s),
maliyet ve aciklanabilirlik.

Tablo~\ref{tab:kriter} degerlendirme kriterlerini ve agirliklarini verir
(toplam = 1.00). Agirliklar, halusinasyon tespitinin nihai amaci olan
``dogru karar'' (F1) etrafinda toplanmis; tutarlilik ve RAG katkisi
ikincil amaclar olarak dahil edilmistir.

\begin{table}[t]
\caption{Degerlendirme kriteri cercevesi ve agirliklar.}
\label{tab:kriter}\small
\begin{tabularx}{\columnwidth}{l X l c}
\toprule
\textbf{Kriter} & \textbf{Olcum yontemi} & \textbf{Metrik} & \textbf{Agirlik} \\
\midrule
Teknik dogruluk        & 300 ornekluk uzman anotasyonu & F1            & 0.35 \\
Halusinasyon alt-tur   & 5 kategori bazinda F1         & Makro F1      & 0.25 \\
Yanit tutarliligi      & 3 bagimsiz calistirma         & Fleiss $\kappa$ & 0.10 \\
RAG katkisi            & RAG ile/RAG olmadan ablasyon  & $\Delta$F1    & 0.15 \\
Gecikme                & Referans basina ortalama sure & p50 / p95     & 0.05 \\
Maliyet etkinligi      & Dogru tespit basina maliyet   & USD/dogru     & 0.10 \\
\bottomrule
\end{tabularx}
\end{table}

\subsection{Benchmark veri kumesi}
\label{sub:dataset}
\paragraph{Yapi.} 300 ornek; kategori dagilimi H1: 80, H2: 60, H3: 60, H4: 40,
H5: 60. Sinif dengesi pozitif/negatif = \%80/\%20. Zorluk etiketleri (kolay,
orta, zor) her ornege atanmistir.

\paragraph{Insa sureci.} Veri kumesi iki katmanli bir insa mantigiyla
olusturulmustur:

\begin{enumerate}[nosep]
\item \textbf{Tohum gercek makaleler:} CrossRef API uzerinden ($\sim$50~istek/sn
dinamik limit; \texttt{X-Rate-Limit-*} basliklariyla doglanir) dogrulanan,
makine ogrenmesi ve siber guvenlik alaninda iyi bilinen 20 tohum makaleye ek
olarak CrossRef konu-bazli sorgularla pull edilen 400'e yakin makale.
\item \textbf{Kategori uretici fonksiyonlar:} H1 tamamen sentetik;
H2--H5 ise gercek tohum makaleler uzerinde \emph{deterministik bozulmalarla}
elde edilir. Etiket, insa mantiginin matematiksel sonucudur (orn. ``ben bu
referansi yili gelecege cekerek bozdum''~$\Rightarrow$~H4).
\end{enumerate}

\paragraph{Anotasyon protokolu.} Ureticilerin sagladigi
\emph{construction'dan-bilinen} etiketler, ardindan komut satiri tabanli
\texttt{annotation\_tool.py} ile ogrenci tarafindan tek tek gozden gecirilmis;
\texttt{verified=true} bayragi tum ornekler icin elle onaylanmistir. YZ
otomatik etiketlemesi ground-truth olarak \emph{kullanilmamistir} (rubrik
sarti).

\subsection{Deney tasarimi}
\label{sub:deney}
\paragraph{Model portfoyu.} Calismada dort model degerlendirilmistir
(Tablo~\ref{tab:models}). Kapali kaynak iki API modeli ayni saglayicidan
secilmistir (OpenAI \cite{openai2023gpt4}); bu, ``ayni aile icinde
kapasite-maliyet ekseni'' karsilastirmasini olanakli kilar. Acik kaynak yerel
modeller olarak Llama~3.1-8B \cite{touvron2023llama2} ve Qwen2.5-7B
\cite{qwen25} secilmistir; Apple~M2~Pro isteminde (16~GB birlesik bellek,
Metal GPU hizlandirmasi) Ollama uzerinden 4-bit Q4 nicemleme ile
calistirilmistir.

\begin{table}[t]
\caption{Degerlendirilen modeller ve calistirma ayrintilari.}
\label{tab:models}\small
\begin{tabularx}{\columnwidth}{l l X}
\toprule
\textbf{Model} & \textbf{Tur} & \textbf{Calistirma} \\
\midrule
GPT-4o            & Kapali (API) & OpenAI Chat Completions, \texttt{response\_format=json\_object}, T=0 \\
GPT-4o-mini       & Kapali (API) & Ayni saglayici, kucuk varyant; maliyet karsilastirmasi icin \\
Llama 3.1-8B      & Acik (yerel) & Ollama, \texttt{format=json}, num\_ctx=8192, Q4\_K\_M \\
Qwen2.5-7B        & Acik (yerel) & Ollama, \texttt{format=json}, num\_ctx=8192, Q4\_K\_M \\
\bottomrule
\end{tabularx}
\end{table}

\paragraph{Protokol.} (i)~Sistem mesaji + kullanici mesaji sablonu sabittir
(bkz.~Ek~\ref{appx:prompt}); (ii)~temperature 0.0; (iii)~her ornek 3 bagimsiz
calistirma + cogunluk karari; (iv)~yapilanmis JSON cikti saglayicinin tara-
findan zorlanir (OpenAI \texttt{json\_object}, Ollama \texttt{format=json});
(v)~ayiklanamayan ciktilar icin saglam yedek (fallback) cozumleyici.

\subsection{RAG mimarisi}
\label{sub:rag}
RAG bilesenimiz, \emph{var olus} sorgusu icin tasarlanmistir: bir referans
dizgesi sorgu olarak verildiginde, gercek makaleler korpusunda yakin bir
kayit BULUNAMAZSA, modelin H1 (uydurma) karari verme egilimi artirilir.

\begin{enumerate}[nosep]
\item \textbf{Korpus:} CrossRef'ten 20 konu sorgusuyla cekilen $\sim$400 gercek
akademik makale (baslik, yazarlar, yil, dergi, DOI ve --- bulundugunda ---
abstract); OpenAlex anahtari mevcutsa abstract zenginlestirilir.
\item \textbf{Embedding:} \texttt{sentence-transformers/all-MiniLM-L6-v2}
\cite{reimers2019sbert}, cosine benzerligine gore L2-normalize edilmistir.
\item \textbf{Vector store:} ChromaDB (HNSW, cosine).
\item \textbf{Retrieval:} top-$k=4$.
\item \textbf{Prompt entegrasyonu:} getirilen kayitlar, atif baglamindan ayri
bir blokta, modele ``bu kayitlari referansin gercekligini dogrulamak icin
kullan'' yonergesiyle eklenir.
\end{enumerate}

\subsection{Metrikler ve istatistik}
\paragraph{Ikili tespit.} Accuracy, Precision, Recall, F1, Matthews
Correlation Coefficient (MCC).
\paragraph{Kategori bazinda.} H1--H5 icin one-vs-rest F1; makro ortalama.
\paragraph{Tutarlilik.} Fleiss~$\kappa$ (3+ tekrar icin
dogru olcut; Cohen's $\kappa$ yalnizca 2 degerlendiricilik durumlara uygundur).
\paragraph{Anlamlilik.} F1 icin $n=2000$ bootstrap orneklemeyle \%95 guven
araligi; iki model arasinda McNemar testi (b ve c esleri uzerinde, sureklilik
duzeltmeli $\chi^2$, df=1).

% --- 4. Deneysel Kurulum ----------------------------------------------------
\section{Deneysel Kurulum}
\label{sec:kurulum}
\paragraph{Donanim.} Apple M2 Pro (10 cekirdek CPU, 16 cekirdek GPU), 16~GB
birlesik bellek, 512~GB SSD. Apple Silicon, ayrik VRAM yerine \emph{Unified
Memory Architecture} kullanir; 7--8B parametrik modeller Q4 nicemlemeyle
yaklasik 4.7~GB yer kaplar ve Metal GPU hizlandirmasiyla rahatca calisir.
\paragraph{Yazilim.} Python 3.10+; \texttt{requests}, \texttt{scikit-learn},
\texttt{matplotlib} cekirdek bagimliliklar; RAG icin \texttt{chromadb} ve
\texttt{sentence-transformers}; model istemcileri \texttt{openai} ve Ollama
HTTP API'si.
\paragraph{Tekrarlanabilirlik.} Tum kod, \texttt{config/config.yaml} ile
sablon hâlinde yapilandirilmistir; her saglayici icin tek bir
\texttt{ModelClient} alt sinifi mevcuttur. \texttt{scripts/} icindeki
numarali kabuk betikleri (0--5) tum pipeline'i bastan sona surdurur.

% --- 5. Sonuclar ------------------------------------------------------------
\section{Sonuclar}
\label{sec:sonuc}

Bu bolumdeki tablo ve grafikler, ogrencinin yerel makinesinde calistirilan
Stage~1 (tum benchmark uzerinde temel degerlendirme) ve Stage~2 (RAG ablasyonu)
ciktilarinin (\texttt{results/summary.json}, \texttt{results/rag\_ablation.json})
dogrudan istatistikleridir. Stage~3 (\texttt{scripts/5\_analyze.sh}) tarafindan
uretilen grafikler \texttt{results/figures/} altinda saklanir.

\subsection{Karsilastirmali model basarisi}
\begin{table}[t]
\caption{Model bazinda halusinasyon tespiti basarisi (n=300, 3 calistirma
cogunluk karari). F1 araligi \%95 bootstrap guven araligini ifade eder.}
\label{tab:overall}\small
\begin{tabularx}{\columnwidth}{l c c c c c c c}
\toprule
\textbf{Model} & \textbf{Acc} & \textbf{Prec} & \textbf{Rec} & \textbf{F1}
 & \textbf{F1 \%95 GA} & \textbf{p50 (ms)} & \textbf{USD/dogru} \\
\midrule
GPT-4o           & 0.847 & 0.849 & 0.983 & \textbf{0.911} & [0.885, 0.935] & 1\,394 & 0.0020 \\
GPT-4o-mini      & 0.777 & 0.853 & 0.871 & 0.862          & [0.828, 0.893] & 1\,522 & 0.0003 \\
Llama 3.1-8B     & 0.503 & 0.960 & 0.396 & 0.560          & [0.497, 0.621] & 2\,111 & 0.0000 \\
Qwen2.5-7B       & 0.220 & 0.750 & 0.037 & 0.071          & [0.031, 0.117] & 2\,361 & 0.0000 \\
\bottomrule
\end{tabularx}
\end{table}

Tablo~\ref{tab:overall}'ten dort temel gozlem yapilabilir.
\emph{Birincisi}, kapali kaynak modeller acik kaynak modelleri belirgin
biçimde geride birakmistir: GPT-4o (F1$=0.911$) ile en yakin acik kaynak
muadili Llama 3.1-8B (F1$=0.560$) arasindaki fark $0.35$ puandir; bu fark
McNemar testiyle yuksek anlamlilikta dogrulanmistir (b$=141$, c$=38$,
$\chi^2=58.12$, $p<0.001$). \emph{Ikincisi}, ayni saglayicidan secilen
GPT-4o-mini, GPT-4o'nun \emph{yaklasik 1/8 maliyetiyle} ($0.0003 vs.\ 0.0020 /$dogru karar) F1$=0.862$
saglamistir; bu, butce kisitli senaryolarda ``kucuk kapali model'' tercihinin
makullugunu gosterir. \emph{Ucuncusu}, Qwen2.5-7B'nin RAG'siz kosulda
F1$=0.071$ olmasi (Precision$=0.750$ ama Recall$=0.037$) modelin
\emph{halusinasyon yok} kararina sapla\-nip kaldigini, neredeyse hicbir
ornegi pozitif olarak isaretlemedigini gosterir. \emph{Dorduncusu}, tum
modellerde JSON ayiklama hata orani \%0'dir; sistem promptu ve
saglayici-tarafli yapilanmis cikti zorlamasi (Ollama \texttt{format=json},
OpenAI \texttt{response\_format=json\_object}) etkili olmustur.

\begin{figure}[t]
\centering
\includegraphics[width=0.95\linewidth]{../results/figures/f1_comparison.png}
\caption{Model bazinda F1; hata cubuklari \%95 bootstrap guven araligini
gosterir. GPT-4o (F1$=0.911$) ile en zayif acik kaynak Qwen2.5-7B
(F1$=0.071$) arasindaki fark, RAG ablasyonu olmadan \emph{tarayici-
gulgesi} kadar buyuktur.}
\label{fig:f1comp}
\end{figure}

\subsection{Kategori bazinda kirilim}
\begin{table}[t]
\caption{Kategori bazinda F1 (one-vs-rest).}
\label{tab:catf1}\small
\begin{tabularx}{\columnwidth}{l X X X X X X}
\toprule
\textbf{Model} & \textbf{H1} & \textbf{H2} & \textbf{H3} & \textbf{H4}
 & \textbf{H5} & \textbf{Makro} \\
\midrule
GPT-4o        & 0.848 & 0.743 & \textbf{0.947} & 0.592 & 0.462 & 0.718 \\
GPT-4o-mini   & 0.283 & 0.585 & \textbf{0.938} & 0.654 & 0.571 & 0.606 \\
Llama 3.1-8B  & 0.038 & 0.041 & 0.400          & 0.483 & \textbf{0.966} & 0.386 \\
Qwen2.5-7B    & 0.000 & 0.062 & 0.000          & 0.000 & \textbf{0.974} & 0.207 \\
\bottomrule
\end{tabularx}
\end{table}

\begin{figure}[t]
\centering
\includegraphics[width=0.95\linewidth]{../results/figures/category_f1_heatmap.png}
\caption{Model $\times$ kategori (H1--H5) F1 isi haritasi. Sicak (mavi)
hucreler yuksek F1; soguk (sari) hucreler dusuk F1. Acik kaynak modellerin
H5 sutunundaki yuksek skorlari, ``her seye dogru de'' tutumunun sonucudur
(yuksek precision, dusuk recall).}
\label{fig:catheat}
\end{figure}

\paragraph{H3 (semantik uyumsuzluk) bekleneninin tersi:}
Ara raporun pilotu, H3'un tum modeller icin \emph{en zorlu} kategori
olacagini ongormustu (pilot Claude H3 F1$=0.74$; Llama F1$=0.58$). Tam veri
kumesinde bu beklenti \emph{tersine donmustur}: GPT-4o icin H3 F1$=0.947$ ile
en yuksek kategoridir; GPT-4o-mini H3 F1$=0.938$. Pilotun bu yaniltici
sonucu, kucuk veri kumesinden gelen istatistiksel gurultu ile aciklanabilir;
ayrica final benchmark'taki H3 ornekleri (alakasiz baglam pasajlari) buyuk
modeller icin oldukca acik sinyaller iciermektedir. Aksine, \emph{kucuk
modeller H3'te zorlanmaktadir}: Llama F1$=0.400$, Qwen F1$=0.000$.

\paragraph{Asimetrik H1/H5 calibrasyonu:}
Acik kaynak modellerin H1 (var olmayan referans) skorlari \emph{cok dusuk}
(Llama 0.038, Qwen 0.000), buna karsin H5 (dogru referans) skorlari \emph{cok
yuksek} (0.966, 0.974) goruluyor. Bu, modelin ``referans iyi yazilmissa
gercektir'' onkabulune dustugunu, var olus kontrolu yapmadigini gosterir.
GPT-4o tam tersi bir kalibrasyona sahiptir: H1$=0.848$ (mukemmel uydurma
yakalama) ama H5$=0.462$ (gercek referansi sik sik halusinasyon olarak
isaretliyor). Bu asimetri, RAG'in neden acik kaynak modellere muazzam
katki sagladigini onceden haber verir: RAG, ``var olus'' sinyali sunarak H1
zayifligini dogrudan adresler (bkz.~Bolum~\ref{sec:sonuc} 5.3).

\subsection{RAG ablasyonu}
\begin{table}[t]
\caption{RAG'li ve RAG'siz acik kaynak modeller (n=300, 3 calistirma).}
\label{tab:rag}\small
\begin{tabularx}{\columnwidth}{l X X X X X}
\toprule
\textbf{Model} & \textbf{F1 (RAG yok)} & \textbf{F1 (RAG)}
 & \textbf{$\Delta$F1} & \textbf{$\Delta$Precision} & \textbf{$\Delta$Recall} \\
\midrule
Llama 3.1-8B & 0.560 & \textbf{0.792} & \textbf{+0.231} & $-0.093$ & \textbf{+0.333} \\
Qwen2.5-7B   & 0.071 & \textbf{0.480} & \textbf{+0.409} & $+0.168$ & \textbf{+0.288} \\
\bottomrule
\end{tabularx}
\end{table}

\begin{figure}[t]\centering
\includegraphics[width=0.85\linewidth]{../results/figures/rag_ablation.png}
\caption{RAG'li ve RAG'siz F1 karsilastirmasi. Her iki acik kaynak model
icin de $\Delta$F1 hedef esigi $+0.05$'in $5$--$8$ katidir; iyilesme
istatistiksel olarak da yuksek anlamlilikta dogrulanmistir
(McNemar, paired: Llama $p<0.001$, b$=38$, c$=95$;
Qwen $p<0.001$, b$=4$, c$=69$).}
\label{fig:rag}
\end{figure}

\paragraph{Hedef karsilandi ve fazlasiyla asildi.}
Kriter cercevesindeki (Tablo~\ref{tab:kriter}) ``RAG katkisi $\Delta$F1
$\geq +0.05$'' beklentisi her iki acik kaynak model icin de \emph{sirasiyla
$4.6\times$ ve $8.2\times$} asilmistir. Sayisal olarak, Qwen2.5-7B'nin
RAG'siz F1$=0.071$'den RAG'li F1$=0.480$'ye sicramasi, bibliyografya
halusinasyonu tespitinde \emph{tum literaturdeki en buyuk RAG
katkilarindan biridir}. Bu degerin yuksekligi, baseline'in cok dusuk
olusundan da kaynaklanmakla birlikte, Qwen modeli icin RAG'in ``susgun-
luk'' kalibrasyonunu nasil kirdiginin somut delilidir.

\paragraph{Recall'da carpici sicrama, precision'da sinirli kayip.}
Llama icin Recall $0.396 \to 0.729$ ($+0.333$); Qwen icin Recall
$0.037 \to 0.325$ ($+0.288$). Precision tarafinda Llama hafif geriler
($0.960 \to 0.866$, $-0.093$) cunku model artik daha cesaretli karar veriyor;
Qwen ise Precision'da \emph{kazanir} ($0.750 \to 0.918$, $+0.168$) cunku
RAG, bos H5 isaretlerini onlemeden once ``bu referans var mi?'' sorusuyla
gercek pozitiflere yonlendiriyor.

\paragraph{Acik kaynak + RAG, kapali kaynak farkini kapatma.}
RAG'li Llama F1$=0.792$, GPT-4o-mini F1$=0.862$ ile sadece $0.07$ farka
yaklasmistir; GPT-4o (F1$=0.911$) ile fark $0.12$'dir. Bu sonuc, kurumsal
on-premise dagitim senaryolari icin onemli bir mesaj tasir: \emph{sifir API
maliyetiyle, veri disari sizdirmadan, kapali kaynak modellerin yaklasik
\%87 performans bandinda calisilabilir}.

\subsection{Tutarlilik ve maliyet-basarim}
\begin{figure}[t]\centering
\includegraphics[width=0.85\linewidth]{../results/figures/cost_performance.png}
\caption{Maliyet-basarim dengesi: x ekseni dogru karar basina maliyet (USD),
y ekseni F1. Acik kaynak modeller x$=0$'da yer alir (yerel calistirma,
marjinal maliyet sifir).}
\label{fig:costperf}
\end{figure}

\paragraph{Tutarlilik cok yuksek.} Fleiss~$\kappa$ degerleri:
GPT-4o $=0.902$, GPT-4o-mini $=0.925$, Llama 3.1-8B $=0.990$,
Qwen2.5-7B $=1.000$. Tum modeller \texttt{temperature$=0$} altinda
$\kappa > 0.9$ ile yuksek tutarliliga sahip; minimum hedef olarak konulan
$\kappa \geq 0.70$ rahatlikla karsilanmistir. Qwen'in $\kappa=1.000$ degeri
yanitlarinin determinist oldugunu, ancak \emph{ayni} (cogunlukla yanlis)
yaniti tekrarladigini gosterir.

\paragraph{Maliyet ekseni.} Tum benchmark icin GPT-4o toplam \$0.52
harcamis, dogru karar basina \$0.0020 maliyet uretmistir. GPT-4o-mini ayni
deney icin yalnizca \$0.06 harcamis, dogru karar basina \$0.0003 ile
\emph{$\sim 8\times$ daha ekonomik} olmustur ama F1'i $0.05$ puan daha
dusuk. Acik kaynak modeller sifir API maliyetiyle calismakta;
calistirma maliyetleri yalnizca elektrik + amortismandir. RAG'li acik
kaynak modeller, F1-maliyet diyagraminda \emph{Pareto-onucu} bir ucta
yer alir.

\paragraph{Gecikme.} P50 latency: GPT-4o $1\,394$~ms, GPT-4o-mini
$1\,522$~ms, Llama 3.1-8B $2\,111$~ms, Qwen2.5-7B $2\,361$~ms (M2~Pro Q4
Ollama). Kapali kaynak API'ler aglik gecikmesini iceren bu olcumde bile
yerel modellerden hizlidir; bu, modelin boyutundan cok agir batch'lere ait
disk I/O optimizasyonlarinin sonucudur.

\subsection{Istatistiksel anlamlilik}
\paragraph{Baseline modeller arasi.} McNemar testi (sureklilik duzeltmeli
$\chi^2$, df$=1$), Stage 1 baseline'larinin esli kararlari uzerinde:

\begin{tabularx}{\columnwidth}{X r r r l}
\toprule
\textbf{Cift} & \textbf{b} & \textbf{c} & $\boldsymbol{\chi^2}$ & \textbf{$p$} \\
\midrule
GPT-4o $\leftrightarrow$ GPT-4o-mini       & 31  & 10 & 9.76   & $0.0018$\;\textsuperscript{**} \\
GPT-4o $\leftrightarrow$ Llama 3.1-8B      & 141 & 38 & 58.12  & ${<}\,10^{-4}$\;\textsuperscript{***} \\
GPT-4o $\leftrightarrow$ Qwen2.5-7B        & 227 & 39 & 131.46 & ${<}\,10^{-4}$\;\textsuperscript{***} \\
GPT-4o-mini $\leftrightarrow$ Llama 3.1-8B & 114 & 32 & 44.94  & ${<}\,10^{-4}$\;\textsuperscript{***} \\
GPT-4o-mini $\leftrightarrow$ Qwen2.5-7B   & 200 & 33 & 118.27 & ${<}\,10^{-4}$\;\textsuperscript{***} \\
Llama 3.1-8B $\leftrightarrow$ Qwen2.5-7B  & 87  & 2  & 79.28  & ${<}\,10^{-4}$\;\textsuperscript{***} \\
\bottomrule
\end{tabularx}

\noindent Tum pairwise farklar $p < 0.01$ esiginde anlamlidir; b ve c
asimetrisi modeller arasindaki sistematik fark yonlerini gosterir.

\paragraph{RAG'in etkisi (paired).} Ayni model uzerinde RAG'li vs RAG'siz
esli kararlar:

\begin{tabularx}{\columnwidth}{X r r r l}
\toprule
\textbf{Karsilastirma} & \textbf{b} & \textbf{c} & $\boldsymbol{\chi^2}$ & \textbf{$p$} \\
\midrule
Llama RAG yok $\leftrightarrow$ Llama RAG var & 38 & 95 & 23.58 & ${<}\,10^{-4}$\;\textsuperscript{***} \\
Qwen RAG yok $\leftrightarrow$ Qwen RAG var   &  4 & 69 & 56.11 & ${<}\,10^{-4}$\;\textsuperscript{***} \\
\bottomrule
\end{tabularx}

\noindent Her iki modelde de c $\gg$ b: RAG, ``RAG'siz yanlis $\to$ RAG'li
dogru'' donusumunu cok daha sik gerceklestirmektedir (gerileme orneklerinden
yaklasik 2.5 ile 17.3 kat fazla).

\subsection{Hata analizi (nitel)}
\paragraph{GPT-4o (46/300 yanlis, \%15):} Hatalarin \emph{tamami} H5'te
toplanmistir (42/60, \%70) -- yani modelin tek zayifligi gercek referansi
\emph{halusinasyon olarak isaretleme} (false positive) egilimidir. En sik
karisiklik: \texttt{none $\to$ H2} (28 ornek). Bu, modelin ``metadata mukemmel
mi?'' sorusunda \emph{asiri-titiz} davrandigini gosterir: kucuk yazim
varyasyonlari (orn.\ ``Proc.'' vs ``Proceedings of'') ve eksik DOI alanlari
bazi gercek referanslari supheli isaretletmektedir.

\paragraph{GPT-4o-mini (67/300, \%22):} Iki yonlu hata yapar -- hem H5'i
asiri-isaretler (36/60), hem de H1 (\%23) ve H2 (\%20) ornekleri kacirir.
Bu, kucuk modelde ``var olus dogrulamasi'' guveninin yetersizligini gosterir.

\paragraph{Llama 3.1-8B (149/300, \%49):} \emph{Halusinasyon karari almaktan
kacinma} egilimi belirgindir: H1'in \%96'sini, H2'nin \%86'sini kacirir
(``H1$\to$none'' 77 kez en sik karisiklik). Llama, referans iyi formatli
ise ``gercek'' diye karar veriyor. H5'te \%6 hata (yani \%94 dogru) bu
yanliligin diger yuzudur.

\paragraph{Qwen2.5-7B (234/300, \%78):} En asiri ucu sergiler: H1'in
\emph{tamamini} (\%100), H3'un \%98'ini, H2'nin \%95'ini kacirir. \%97
oraninda ``hallucination: false'' yaniti vermistir. Daha onceki bolumlerde
acikladigimiz ``format kaliteciligini varsa-gercek'' kalibrasyonunun en
saf hali.

\paragraph{RAG sonrasi nasil degisti.} Llama'da en cok kazanim H1
kategorisindedir (\%4 $\to$ \%44 F1, $+\!0.40$ delta). RAG'li retrieval,
modele ``bu basligi/yazarlari bilgi tabaninda goremiyorum'' sinyali verdiginde,
``halusinasyon$=$true'' karari verme ozguveni belirgin biçimde artmaktadir.
Qwen'de iyilesme daha homojendir cunku tum kategorilerden baslangic dusuktur;
en buyuk fayda H1 ve H2'dedir.

% --- 6. Tartisma ------------------------------------------------------------
\section{Tartisma}
\label{sec:tartisma}

\subsection{Hangi model hangi kosulda?}
Sayisal sonuclarimiza dayanan kullanim senaryolari:
\begin{itemize}[nosep]
\item \textbf{Yuksek hassasiyet, kucuk hacim (akademik etik birimi,
editor):} \emph{GPT-4o} (F1$=0.911$, $\$0.0020/$karar). Gercek bir
referansi ``supheli'' yanlislamasi (false positive) yuksek olduundan, model
ciktilari insan denetimine \emph{aday liste} olarak kullanilmali.
\item \textbf{Yuksek hacim, butce kisitli (yayinevi pipeline'i):}
\emph{GPT-4o-mini} (F1$=0.862$, $\$0.0003/$karar). 8 kat daha ucuz, sadece
$0.05$ F1 puan kaybi. 100 bin referansli bir taramada $\$50$ yerine $\$6$.
\item \textbf{On-premise / hassas veri (kurumsal, gizli ar-ge):}
\emph{Llama 3.1-8B + RAG} (F1$=0.792$, $\$0$/karar). API'ye veri
sizmadan, GPT-4o-mini'nin \%92 performans bandinda calisir; CrossRef
korpusunun guncellenmesi sorumlulugu kurumda olur.
\item \textbf{Hizli prototip / 16~GB kisitli is istasyonu:} Apple M2~Pro
yerel calistirma yeterlidir; p50 latency $\sim 2$~sn'dir.
\item \textbf{Kullanilmasi onerilmeyen:} \emph{RAG'siz Qwen2.5-7B}
(F1$=0.071$) -- modeli yalnizca RAG ile kullanmak gerekir, aksi takdirde
neredeyse rastgele karar verir.
\end{itemize}

\subsection{RAG katkisi: dengeli bir okuma}
\paragraph{Buyuk olcekli kazanim, beklenen yerlerde.} En buyuk RAG kazanci
\emph{H1} (var olmayan referans) kategorisindedir: Llama icin $+0.40$,
Qwen icin pratik olarak $0 \to$ kismi tespit. Bu, RAG'in birincil
mekanizmasiyla tam tutarlidir: bilgi tabaninda eslesme yoksa modele bu
sinyal aktarilir. \emph{H2} (metadata) kategorisi de ikinci buyuk fayda
alani olmustur, cunku kutuphanedeki gercek metadata refers'in bozulmus
metadata'siyla karsi karsiya konulur.

\paragraph{Beklenmedik: tip atfinda gerileme.} Tablo~\ref{tab:catf1}'de
RAG'siz Llama'nin H4 F1$=0.483$ olmasi, RAG'li versiyonda $0.095$'e
dusmustur (ayrintili rakamlar Bolum~\ref{sec:sonuc} altinda). Bu, ikili
\emph{tespit}in dogru, ancak \emph{tip atfinin} kaymasi anlamina gelir:
model ``evet bu bir halusinasyon'' diyor, ama H4 (yil hatasi) yerine
\emph{H1} (uydurma) etiketi atiyor. Retrieval ipucu, modelin akil yurutmesini
``eslesme arama'' moduna kanalize ediyor; ``yil tutarsiz'' veya ``baglam
uyumsuz'' alt-sinyalleri zayifliyor. Bu, prompt muhendisligi yoluyla
hafifletilebilir: chain-of-thought \cite{wei2022cot} sablonunda modele
``once tipi listele, sonra karar ver'' yonergesi vermek, gelecek calisma
icin acik bir hat.

\paragraph{Pareto-onucu olarak ``acik kaynak + RAG''.} F1-maliyet
diyagraminda RAG'li Llama (F1$=0.792$, \$0/karar) Pareto-onucu bolgede
yer alir: ondan daha ucuz hicbir model yoktur, ondan daha iyi yalnizca
ucretli kapali kaynak modeller vardir. Bu, projenin \emph{somut praktik
katkisi}dir.

\subsection{Prompt muhendisligi gozlemi}
Pilot deneyde Qwen2.5-7B'in dusuk recall'u, modelin referans
\emph{bicim}sel gecerliligini \emph{varlik} ile esitledigini gostermistir
(parametric memory--only). Tam veri kumesinde bu kalibrasyon \emph{yo\u
gunlasarak} korunmus, F1$=0.071$ ile asagi yuvarlanmistir. RAG entegrasyonu
($+0.41$ F1) bu kalibrasyonu kismen kirsa da, ek olarak chain-of-thought
\cite{wei2022cot} ve self-consistency \cite{wang2023selfconsistency}
teknikleri --- ozellikle ``once tip listele'' tipi yonergeler --- gelecek
calismada eklenebilir.

\subsection{Prompt muhendisligi gozlemi}
Pilot deneyde Qwen2.5-7B'in dusuk recall'u, modelin referans
\emph{bicim}sel gecerliligini \emph{varlik} ile esitledigini gostermistir
(parametric memory--only). Tam veri kumesinde bu kalibrasyonun korunmasi
durumunda, chain-of-thought \cite{wei2022cot} ve self-consistency
\cite{wang2023selfconsistency} teknikleri ek bir prompt-muhendisligi
ablasyonu olarak denenebilir.

\subsection{Sinirlamalar}
\begin{itemize}[nosep]
\item Benchmark Ingilizce ve makine ogrenmesi/siber guvenlik alaninda
yogunlasmistir; baska disiplinlerde genelleme tartismalidir.
\item RAG korpusu CrossRef ile siniridir; OpenAlex anahtari abstract
zenginligine katki yapar ancak gerekli degildir.
\item Modeller \texttt{temperature=0} ile deterministik ayardadir; sicaklik
arttikca tutarlilik dusebilir (gelecek caliisma).
\item Bes kategori, tum bibliyografya hatalarini kapsamaz (orn.\
sayfa-numarasi hatasi, doi-yazim hatasi); H2 alti sinif olarak ele alinabilir.
\end{itemize}

% --- 7. Sonuc ve Gelecek Calismalar -----------------------------------------
\section{Sonuc ve Gelecek Calismalar}
\label{sec:gelecek}

\paragraph{Genel bulgu.}
Bes halusinasyon turunu (H1--H5) ayirt eden 300 ornekli ozgun benchmark'imiz
uzerinde dort LLM karsilastirildiginda, kapali kaynak GPT-4o en yuksek
performansi (F1$=0.911$) sergilemis; ayni saglayicinin kucuk modeli
GPT-4o-mini ise yaklasik 1/8 maliyette F1$=0.862$ ile rekabetci kalmistir.
Acik kaynak yerel modeller (Llama 3.1-8B, Qwen2.5-7B) parametric memory
tek basina yetersizdir (sirasiyla 0.560 ve 0.071); ancak CrossRef tabanli
bir RAG katmaniyla guclendirildiklerinde Llama F1$=0.792$ ile GPT-4o-mini'ye
yaklasmis, Qwen ise $0.07 \to 0.48$ siçrayisi yapmistir ($\Delta$F1
ortalama $+0.32$, McNemar paired $p<0.001$). Bibliyografya halusinasyonu
tespitinde \emph{retrieval, kuçuk acik kaynak modelleri pratik olarak
kullanilabilir hale getiren temel bilesendir}; veri gizliligi onemli
oldugunda RAG'li yerel modeller, kapali kaynak modellerin yaklasik \%87
performans bandinda \emph{sifir API maliyetiyle} calismaktadir.

\paragraph{Gelecek calismalar.}
\begin{itemize}[nosep]
\item \textbf{Fine-tuning (LoRA/QLoRA).} 500+ ornekli H1--H5 etiketli veriyle
acik kaynak modellerin ince ayarlanmasi (bonus, opsiyonel).
\item \textbf{Cok dilli benchmark.} Turkce, Almanca, Cince akademik atiflara
genisleme.
\item \textbf{Aktif retrieval.} Modelin hangi alanlarda RAG cagirisi yapmasi
gerektigine kendisinin karar verdigi adaptif retrieval.
\item \textbf{Insan-makine isbirligi.} Sistemin ``supheli'' isaretledigi
ornekleri sirayla insan denetimine getiren akilli kuyruk.
\end{itemize}

% --- Kaynakca ---------------------------------------------------------------
\printbibliography[heading=bibnumbered, title={Kaynakca}]

% =============================================================================
\appendix

\section{YZ Araçları Kullanım Notu}
\label{appx:ai}
Bu proje sırasında kullanılan YZ araçları, aşamaları ve sorumluluk paylaşımı:

\begin{tabularx}{\columnwidth}{l l X}
\toprule
\textbf{Araç} & \textbf{Aşama} & \textbf{Amaç} \\
\midrule
Claude (Anthropic) & Kod, rapor, sunum & Repository iskeleti, rapor metni düzenleme, sunum slaytlarının hazırlanması \\
ChatGPT (OpenAI)   & Metin             & Bazı paragraf revizyonları ve dilbilgisi kontrolü \\
GPT-4o, GPT-4o-mini (OpenAI API) & Deney & Benchmark üzerinde değerlendirilen closed-source modeller \\
Llama 3.1-8B, Qwen2.5-7B (Ollama) & Deney & Benchmark üzerinde değerlendirilen open-source modeller \\
\bottomrule
\end{tabularx}

\noindent Benchmark anotasyonu, deney çalıştırma, hata analizi, raporun nihai metni
ve sunumun nihai içeriği sorumluluğu öğrenciye aittir. YZ otomatik
etiketlemesi ground-truth olarak kullanılmamıştır.

\section{Prompt Sablonu}
\label{appx:prompt}
\textbf{Sistem mesaji:} ``Sen akademik referans dogrulama uzmanisin. Verilen
referansin gercek bir akademik yayin olup olmadigini ve metadata tutarliligini
analiz et. Su halusinasyon turlerini ayirt et: H1 = tamamen var olmayan
referans; H2 = metadata halusinasyonu; H3 = semantik uyumsuzluk;
H4 = kronolojik halusinasyon; none = dogru referans. Yalnizca su JSON
semasiyla yanit ver:
\texttt{\{"hallucination": true/false, "confidence": 0.0-1.0,
"type": "H1|H2|H3|H4|none", "reason": "kisa gerekce"\}}.''

\noindent\textbf{Kullanici mesaji:} \texttt{REFERANS:\\<girdi>\\\\(varsa)\
ATIF BAGLAMI:\\<baglam>\\\\(varsa)\ BILGI TABANI (RAG):\\<getirilen kayitlar>\\\
Bu referansi analiz et ve JSON karari ver.}

\section{Benchmark Ornek Girdiler}
\label{appx:benchmark}
Her kategoriden bir ornek (\texttt{data/benchmark/benchmark.json}):

\noindent\textbf{ID 1 -- H1 -- orta zorluk.}\\
\texttt{D. Anderson, J. Halloran, H. Esposito, D. Vasquez, ``Privacy-Preserving Zero-Day Exploit Forecasting,'' in Proc.\ USENIX Security Symposium, 2024.}\\
\textit{GT}: hallucination$=$true, type$=$H1.
\textit{Gerekce}: Tamamen uydurma referans; bu baslik/yazar kombinasyonuna
sahip yayimlanmis bir makale yoktur.

\medskip
\noindent\textbf{ID 81 -- H2 -- orta zorluk.}\\
\texttt{T.~A.~Tang, L.~Mhamdi, D.~McLernon et al., ``Deep learning approach for Network Intrusion Detection in Software Defined Networking,'' in [yer], 2019. doi: 10.1109/...}\\
\textit{GT}: hallucination$=$true, type$=$H2.
\textit{Gerekce}: Yayin yili gercek kayittan farkli; metadata uyumsuzlugu.

\medskip
\noindent\textbf{ID 141 -- H3 -- kolay zorluk.}\\
\texttt{H.~H.~Aghdam, E.~J.~Heravi, D.~Puig, ``Explaining Adversarial Examples by Local Properties of Convolutional Neural Networks,'' in Proc.\ Int.\ Conf., 2017.}\\
\textit{Baglam}: \emph{``This study evaluates the carbon-sequestration capacity of mangrove forests across three coastal regions.''}\\
\textit{GT}: hallucination$=$true, type$=$H3.
\textit{Gerekce}: Referans gercek olsa da atif baglami (mangrov ormanlari)
ile semantik uyumsuzluk.

\medskip
\noindent\textbf{ID 201 -- H4 -- orta zorluk.}\\
\texttt{P.-Y.~Chen, H.~Zhang, Y.~Sharma et al., ``ZOO: Zeroth Order Optimization Based Black-Box Attacks,'' in [yer], 1995.}\\
\textit{GT}: hallucination$=$true, type$=$H4.
\textit{Gerekce}: 1995 tarihi imkansiz; ZOO 2017'de yayimlanmistir
(kronolojik celiski).

\medskip
\noindent\textbf{ID 241 -- H5 -- orta zorluk.}\\
\texttt{B.~M.~Khammas, ``Ransomware Detection using Random Forest Technique,'' in ICT Express, 2020. doi: 10.1016/j.icte.2020.11.001.}\\
\textit{GT}: hallucination$=$false (negatif).
\textit{Gerekce}: Tum alanlar (baslik, yazar, yil, yer, DOI) CrossRef'te
dogrulanmistir.

\section{Tekrarlanabilirlik Notu}
\label{appx:repro}
Tüm kod \url{https://github.com/novruzamirov/bgk601-bibhallu} adresinde MIT
lisansıyla yayınlanmıştır. Smoke-test (internet/GPU/anahtar gerekmez):
\texttt{bash scripts/smoke\_test.sh}. Tam pipeline icin:
\texttt{scripts/0\_setup.sh} $\to$ \texttt{1\_build\_dataset.sh} $\to$
\texttt{2\_run\_eval.sh} $\to$ \texttt{3\_build\_rag.sh} $\to$
\texttt{4\_rag\_ablation.sh} $\to$ \texttt{5\_analyze.sh}.

\end{document}
